\section{Introduction}
\label{sec:intro}

Modern stock exchanges such as Nasdaq and NYSE run two trading systems~\cite{harris2003}. During regular hours, traders submit orders to a continuous limit order book (CLOB), where a buy order and a sell order execute as soon as their prices meet. Traders can also submit orders to a scheduled auction, or call market, which holds orders until a fixed time and then executes them together at a single price. The closing auction sets the official closing price of each stock, against which portfolios are benchmarked, indexes are computed, and derivatives are settled; closing auctions account for 7.5\% of daily volume in U.S.\ stocks~\cite{bogousslavsky2023}. The same mechanism opens the market the next morning.

Most activity still occurs in the CLOB, and a CLOB processes orders in the order they arrive. A difference of microseconds therefore decides who trades first. High-frequency traders (HFTs) invest in fast hardware, private data links, and co-location to win these races. Aquilina, Budish, and O'Neill measure them directly on exchange message data: races account for about 20\% of trading volume in FTSE~100 stocks and the sums at stake are on the order of \$5~billion a year across global equity markets~\cite{aquilina2022}. HFTs also narrow spreads and add liquidity, but liquidity that can be withdrawn faster than other traders can react is of limited use to them.

\paragraph{Frequent batch auctions.}
In an opening or closing auction, reaching the exchange a few microseconds earlier brings no advantage, because orders are not matched one by one. Budish, Cramton, and Shim propose to run such an auction repeatedly through the day~\cite{budish2015,budish2014}. In a frequent batch auction (FBA), the exchange collects orders over a short interval (\eg one second) and then clears the whole batch at one price. An order that arrives at the start of the batch is treated the same as one that arrives at the end: there is price priority and no time priority. Traders compete on price and no longer on speed.

\paragraph{The transparency gap.}
An FBA gives up something a CLOB has. In a CLOB the book is visible, so once an order reaches the matching engine there is no doubt about why it executed, did not execute, or executed at a given price. In a batch auction, no participant sees the set of orders that determines the clearing price, and after the close participants see only the outcome. The orders that failed to execute stay hidden, although the clearing price depended on them. A trader observes whether their own order executed and at what price, and cannot check that the exchange included every eligible order, excluded ineligible ones, chose the price the rule prescribes, and allocated trades correctly. The exchange cannot repair this by publishing the book, since unexecuted orders reveal trading interest that their owners have reason to keep private.

\paragraph{A proof instead of the book.}
A zero-knowledge proof~\cite{gmw1986} shows that a computation was performed correctly without revealing its private inputs. Applied to a batch auction, the exchange proves that the published clearing price and volumes follow from the submitted orders and reveals none of them. Custodial cryptocurrency exchanges already publish proofs of this kind about their hidden balance sheets~\cite{provisions2015}; we apply the approach to the matching engine.

We develop a zero-knowledge frequent batch auction (zk-FBA). The exchange commits to the order book as orders arrive, returns to each trader a receipt that the order was recorded, and publishes at the close a clearing receipt with a succinct proof (a zk-SNARK) that the receipt is the outcome the auction rule prescribes for the committed book. The exchange must produce this proof once per batch, which is the binding constraint on the design. A general-purpose proof system compiles the clearing rule into a circuit and pays for the generality at every batch. We instead write the rule directly as polynomial identities over tick-indexed columns: a custom arithmetization~\cite{plonkbook} for a polynomial interactive oracle proof (PIOP)~\cite{dark2020}, in the style of PLONK~\cite{plonk2019} and with KZG commitments~\cite{kzg2010}.

An FBA has no canonical specification, since it is a market-design proposal and not an operational standard~\cite{budish2014}. We fix the details by adopting the rule of a deployed auction, the Nasdaq Closing Cross, and assume only that it runs frequently.

\subsection{Contributions}
\label{sec:contributions}

\begin{enumerate}
  \item \textbf{A custom arithmetization of the clearing rule.} We encode a closing-cross clearing rule, with limit, market, and imbalance-only orders and its three tie-breaks, as 33 polynomial constraints over 88 columns indexed by price tick, combined into one quotient. We prove that the constraints admit exactly one clearing receipt for a given book (\S\ref{sec:protocol}).
  \item \textbf{Lookup range checks.} Every inequality in the rule is a range check. We replace one committed column per bit with 8-bit limbs and a single logarithmic-derivative lookup, which reduces the proof for 32-bit volumes from 352 commitments and 22.9\,KB to 89 commitments and 6.1\,KB (\S\ref{sec:protocol}).
  \item \textbf{A committed, updatable book.} We commit to the book with a homomorphic polynomial commitment, so that the exchange updates the commitment for each order, returns a receipt to the trader, and hides the side of the order in the public log. A second proof at the close shows that every logged update is a valid order (\S\ref{sec:book}).
  \item \textbf{Two implementations.} We implement the prover and verifier in Rust on arkworks, and independently write the clearing statement as a Noir circuit proved with Barretenberg. For a book of 1,000 price ticks, we generate the clearing proof in 1.34\,s on one core of a laptop and 0.37\,s on eight threads; the proof is 6,124 bytes and verifies in 29\,ms. This is 5.5 times faster than the circuit, with a proof 2.4 times smaller (\S\ref{sec:performance}).
  \item \textbf{Order flow from a real market.} No batch-auction venue publishes order-level data. We replay Nasdaq TotalView-ITCH order flow as one-second batches, measure how often each tie-break is reached, and build a calibrated generator of benchmark books. Both implementations give the same receipt on all 32 real and 25 synthetic books (\S\ref{sec:performance}).
\end{enumerate}
